% MT65091C
% Vergiftungen, allgemein
% 14.0
% 2013-11-30
\label{m:am-vergiftung}
\begin{enumerate}
  \item Selbstschutz
  \item Beendigung des Giftkontaktes
  \item Evtl. Dekontamination bei
  Gefahrstoffen\footnote{z.\,B. Insektizide, Lösungsmittel
  in der chemischen Industrie, \ldots}
  \item Sicherung der Vitalfunktionen
  \item Entscheidung, ob es sich um einen akut vital
  bedrohten Patienten handelt
  \item Quellenforschung (Mistkübel
  durchsuchen)
  \item Nüchtern lassen
  \item NICHT absichtlich Erbrechen
  herbeiführen
  \item Asservierung (Packungen, Erbrochenes in
  Nierentasse mitnehmen!)
  \item Anamnese, Besonderheiten:
  \begin{multicols}{2}
  \begin{enumerate}\raggedbottom
    \item \EE{Wer}? (Alter, KG, Vorerkrankungen, \ldots
    $\rightarrow$ Risikoabschätzung!)
    \item \EE{Was}? (Art des Giftes $\rightarrow$ evt.
    zu erwartende Symptome)
    \item \EE{Wann}? (Zeit seit Einnahme)
    \item \EE{Wie}?
    (inhalativ, oral, perkutan, intravenös etc.)
    \item \EE{Wieviel}? (Dosisabschätzung)
    \item \EE{Warum}? (Ursachenforschung:
    Selbst\-mord\-ver\-such, Unfall, Fremdverschulden,
    Vorsatz?)
  \end{enumerate}
  \end{multicols}
  \item Ggfs. \EE{Expertenrat einholen} und Entscheidung
  bezüglich des weiteren Vorgehens:
  \begin{multicols}{2}
  \begin{enumerate}\raggedbottom
    \item Gegenmittel verfügbar?
    \item \enquote{Normales} Spitalsbett?
    \item Intensivüberwachung?
    \item Intensivüberwachung mit Entgiftungseinheit?
    \item Dialyse?
    \item Druckkammer?
    \item Sonstige Spezialbehandlungseinheit?
  \end{enumerate}
  \end{multicols}
\end{enumerate}